\documentclass{article}
\usepackage{graphicx} % Required for inserting images
\usepackage{amsmath}
\usepackage{amsfonts}
\usepackage{listings}
\usepackage{booktabs}
\usepackage{float}
\usepackage{hyperref} % should always go last
% \usepackage{enumitem}

\lstset{
  basicstyle=\ttfamily\footnotesize,
  frame=single,
  captionpos=t,
  breaklines=true,
}

\newcommand{\includesource}[1]{%
  \lstinputlisting[caption={\texttt{#1}}]{#1}%
}

\title{Project Report - Group 8}
\author{Hans Bouwman, Joel de Haan}

\begin{document}

\maketitle


\tableofcontents

\newpage




There are two obvious low-hanging fruits when it comes to improving performance. Consider the following loop in the solver:

\begin{lstlisting}
for (int ix = 1; ix < (nx-1); ix++)
{
    for (int iy = 1; iy < (ny-1); iy++)
    {
        v[nx*iy+ix] = vp[nx*iy+ix];
        w += fabs(v[nx*iy+ix]);
    }
}
\end{lstlisting}

Firstly, the loops should be swapped to make access contiguous, as this improves cache coherence. Secondly, \texttt{vp} is not read before being written to again in the next iteration.
As such, copying the array is completely superfluous and we can simply swap the pointers. As such we suggest the following code change:

\begin{lstlisting}
swap(v, vp);
for (int iy = 1; iy < (ny-1); iy++)
{
    for (int ix = 1; ix < (nx-1); ix++)
    {
        w += fabs(v[nx*iy+ix]);
    }
}
\end{lstlisting}


\end{document}
